\section{Away steps: certificates and removal of active atoms}
Let $Q=\operatorname{conv}(\mathcal A)$ for a finite atom set. Maintain
$x=\sum_{a\in S}\lambda_a a$ with positive weights summing to one.
Besides the FW atom $s$, choose $v\in\arg\max_{a\in S}\inner{\nabla f(x)}{a}$.
Use the larger descent gap of $s-x$ and $x-v$; ties choose FW. Denote the
chosen direction by $d$ and its gap by $g=-\inner{\nabla f(x)}d$.
Its feasible cap is $\alpha_{\max}=1$ for FW and
$\lambda_v/(1-\lambda_v)$ for an away step. A singleton active set takes
a FW step. The stopping certificate always uses the \emph{full FW gap},
not the selected direction's gap. In particular $g\ge G\ge U$.

This is classical away-step FW with standard quadratic
backtracking~\cite{lacostejaggi,pedregosa}. Start from $L_{-1}=M>0$;
each update halves the previous estimate, doubles on rejection, and
uses
\[
 \alpha=\min\{\alpha_{\max},g/(L\|d\|^2)\}
\]
in~\eqref{eq:accept}, with $G$ there replaced by $g$.
An away update at its cap is a \emph{drop step}; all other updates are
called good. Searching the active set incurs additional work, separate
from one full LMO per iterate.
The bounds in this section use exact arithmetic and exact full and
active-set linear oracles.

\begin{lemma}[Local acceptance and counting drops]\label{lem:away}
The quadratic test passes for $L\ge H_0+H_1\delta$, where
$\delta=f(x)-f_Q^*$. This also holds when $\alpha_{\max}>1$.
If $m$ good updates occur, the total number of updates is at most
$2m+|S_0|-1$.
\end{lemma}
\begin{proof}
On the feasible segment $[0,\alpha_{\max}]$, the same initial-descent
argument as in Lemma~\ref{lem:progress} applies with the chosen $g$.
More explicitly, while the function is at most $f(x)$ its directional
second derivative is bounded by $(H_0+H_1\delta)\|d\|^2$.
The resulting quadratic upper bound is strictly below $f(x)$ for
$0<t\le\min\{\alpha_{\max},g/(L\|d\|^2)\}$, preventing a first
return to the initial level. Thus the bound holds through the trial.
Every drop removes an atom, whereas a FW step adds at most one and is
good. Active-set cardinality is positive, so drops number at most the
number of FW steps plus $|S_0|-1$.
\end{proof}

\begin{theorem}[Convex certificate bound with away steps]\label{thm:awayconvex}
With $J,P$ as in Theorem~\ref{thm:ordinary}, away-step backtracking
reaches $U_N\le\varepsilon$ in at most
\[
 2\left[P+\left\lceil(6+16H_1D^2)J+32H_0D^2/\varepsilon\right\rceil\right]
 +|S_0|-1
\]
updates. It uses $N+1$ gradients and full LMOs, and, for actual $N\ge1$,
\[
 N_f\le1+2N+\left\lceil\log_2\max\{1,2(H_0+H_1U_0)/M\}\right\rceil.
\]
\end{theorem}
\begin{proof}
Use the same phase caps and integer potential as in
Theorem~\ref{thm:ordinary}, counting lag steps across \emph{all} updates.
On a good away step the quadratic minimizer is not clipped; on a full
FW step acceptance decreases $f$ by at least $g/2$. Hence a good update
with $L\le C_j$ decreases $f$ by at least
$\min\{b_j/4,b_j^2/(8C_jD^2)\}$. Count only these good updates using
the phase budget; drops also decrease $f$ and cannot increase that
budget. The resulting bound on good updates is the ordinary bound.
Apply Lemma~\ref{lem:away}. The accepted-curvature envelope and exact
telescoping trial identity hold across drops as well, proving the
function-call bound.
\end{proof}
The known-$H$ AFW quadratic rule, with $L=H_0+H_1U_k$ and the same
cap, similarly inherits Theorem~\ref{thm:main}'s bound on good updates.
Its total updates are at most twice that bound plus $|S_0|-1$, and it
uses $N+1$ function calls. The stronger theorem below concerns ordinary
backtracking, whose accepted estimates can follow the true local error.

\subsection{A logarithmic certificate bound under a geometric error bound}
Assume the selected gaps satisfy the standard geometric error bound
\begin{equation}\label{eq:awaygeometry}
 g_k^2\ge 2\kappa\bigl(f(x_k)-f_Q^*\bigr),\qquad\kappa>0.
\end{equation}
Strong convexity with parameter $\mu>0$ is one sufficient condition:
for a finite polytope, one may use
$\kappa=\mu W^2/4$, where $W$ is its pyramidal width. The pairwise
gap squared is at least $2\mu W^2\delta$ and the selected gap is at
least half that pairwise gap~\cite{lacostejaggi,pedregosa}.
No numerical width estimate fitted to a trajectory is used here.
Strong convexity is not required separately when~\eqref{eq:awaygeometry}
is available; geometric error-bound rates are established
results~\cite{wirth,restartfw}.

\begin{proposition}[A standard quadratic-growth sufficient condition]
\label{prop:qg}
Let $X^*=\arg\min_Q f$ and suppose that, on the initial sublevel set,
\[
 f(x)-f_Q^*\ge\frac\gamma2\operatorname{dist}(x,X^*)^2,
 \qquad\gamma>0.
\]
For a polytope of pyramidal width $W$, condition~\eqref{eq:awaygeometry}
holds along monotone iterates with $\kappa=\gamma W^2/16$.
No unique minimizer or positive strong-convexity parameter is required.
\end{proposition}
\begin{proof}
Let $p$ be the pairwise gap for the current active set, and take a
nearest $x^*\in X^*$ when $\delta>0$. The standard polytope scaling
inequality and convexity give
\[
 p\ge W\frac{\inner{\nabla f(x)}{x-x^*}}{\norm{x-x^*}}
   \ge W\frac\delta{\operatorname{dist}(x,X^*)}
   \ge W\sqrt{\gamma\delta/2}.
\]
The selected FW/away gap is at least $p/2$. Squaring proves
$g^2\ge\gamma W^2\delta/8=2\kappa\delta$; zero error is immediate.
Monotone iterates stay in the sublevel set where growth is assumed.
\end{proof}
This is the known scaling/Holderian-error-bound implication at exponent
$1/2$~\cite{restartfw}, Lemmas 3.4 and 3.6, with constants translated
to our selected-gap convention. It is not a new geometric result.
The direct strong-convexity implication above has a better constant.
For example, $f(x)=\cosh(x_1)-1$ on a two-dimensional $\ell_1$ ball
satisfies quadratic growth with $\gamma=1$ and a non-singleton optimal
set, while its Hessian has a zero eigenvalue.

\begin{lemma}[From primal error to the stopping certificate]\label{lem:awaybridge}
Under~\eqref{eq:curvature}, every feasible point with primal error
$\delta=f(x)-f_Q^*$ satisfies
\[
 U\le G\le\widehat B(\delta):=
 \begin{cases}
 \sqrt{2C(\delta)\delta},&2\delta\le C(\delta),\\
 \delta+C(\delta)/2,&2\delta>C(\delta),
 \end{cases}
 \qquad C(\delta)=D^2(H_0+H_1\delta).
\]
This is at most the earlier conservative envelope
$B(\delta)=\max\{2\delta,\sqrt{2C(\delta)\delta}\}$.
For $D>0$, put $a=D^2H_0$, $b=D^2H_1$.
Then $\delta\le\eta_\varepsilon$ implies $U\le\varepsilon$, where
\[
 \eta_\varepsilon=
 \begin{cases}
 (\varepsilon-a/2)/(1+b/2),& b<2\text{ and }\varepsilon(2-b)\ge2a,\\[2pt]
 \varepsilon^2/(a+\sqrt{a^2+2b\varepsilon^2}),&\text{otherwise}.
 \end{cases}
\]
In particular $H_0=H_1=0$ gives $\eta_\varepsilon=\varepsilon$.
Neither strong convexity nor uniqueness of the minimizer is needed.
\end{lemma}
\begin{proof}
Use the local quadratic model on the hypothetical step
$t=\min\{1,G/C(\delta)\}$, with a full step if $C(\delta)=0$.
The initial-descent argument in Lemma~\ref{lem:progress} proves that
the model is valid through this step; $\norm{s-x}\le D$.
Feasibility and the exact model decrease give
\[
 \delta\ge\Phi_C(G):=
 \begin{cases}G^2/(2C),&G\le C,\ C>0,\\
 G-C/2,&G>C\text{ or }C=0.
 \end{cases}
\]
Inverting this increasing function of $G$ gives $\widehat B$.
If $2\delta\le C$, the square-root branch is unchanged; otherwise
$\delta+C/2<2\delta$ and exceeds $\sqrt{2C\delta}$.
Solving $\widehat B(\eta)=\varepsilon$ gives the stated threshold.
The full-step root applies exactly when $b<2$ and
$\varepsilon(2-b)\ge2a$; otherwise rationalization gives the
quadratic root. This threshold is at least the earlier conservative
$\min\{\varepsilon/2,\varepsilon^2/(a+\sqrt{a^2+2b\varepsilon^2})\}$,
with the zero-curvature convention for that expression.
\end{proof}
The two-branch inversion already appears in~\cite[Eq.~(15)]{lacoste2016}.
Here the feasible-descent argument uses a source-dependent local
curvature bound and retains the exact full-step decrease.
For $H_0=0<H_1$, the threshold is
$\varepsilon/(1+b/2)$ when $b<2$, and
$\varepsilon/\sqrt{2b}$ when $b\ge2$. For fixed $H_0>0$ and small
$\varepsilon$, it is still of order $\varepsilon^2$.

The square-root branch cannot in general be replaced by a linear
primal-to-gap bound. For $f(x)=x^2/2$ on $[-1,1]$, at $x=t>0$,
\[
 \delta=t^2/2,\qquad G=t(1+t),\qquad G/\delta=2(1+t)/t\longrightarrow\infty.
\]
At the first observation the scalar certificate also equals $G$.
This example has strong convexity one. It limits pointwise conversion
for this certificate, not the iteration complexity of a method or
the benefit of an accumulated history. In fact, $\widehat B$ is attained
on both branches when $H_1=0$: for $Q=[-D/2,D/2]$, a source at $D/2$
and a quadratic of curvature $H_0$ centred inside $Q$ give the
square-root branch; centres to the left of $-D/2$ give the full-step
branch. Thus this refinement is sharp for the quadratic descent
inequality; no optimal-envelope claim is made for general $H_1>0$.

\begin{theorem}[Phase bound under a geometric error bound]\label{thm:awaystrong}
Assume additionally $H_0+H_1>0$ and~\eqref{eq:awaygeometry}. For
$0<\varepsilon<U_0$, set
\[
 J_*=\left\lceil\log_2(U_0/\eta_\varepsilon)\right\rceil,\qquad
 P_*=\left\lceil\log_2\max\{1,M/[2(H_0+H_1U_0)]\}\right\rceil.
\]
The same algorithm reaches $U_N\le\varepsilon$ within
\[
 N\le2\left\lceil P_*+
       (6+4H_0D^2/\kappa)J_*+8H_1D^2U_0/\kappa\right\rceil+|S_0|-1.
\]
The preceding function-call bound remains valid. The algorithm needs
none of $H_0,H_1,\kappa,f_Q^*$ as inputs. For fixed problem and
initialization, the accuracy dependence is logarithmic.
\end{theorem}
\begin{proof}
Analyze phases of the \emph{true} error, with $b_j=U_0/2^j$, stopping
at $\eta_\varepsilon$. These phases are for analysis only. Use positive
caps $C_j=2(H_0+H_1b_j)$, which decrease by at most a factor two.
The initial potential is at most $P_*$ in virtual phase zero. Allowing
initially skipped phases, its total increases number at most $J_*-1$;
at most $P_*+J_*$ lag steps occur.

On a controlled good step, acceptance and~\eqref{eq:awaygeometry} give
\[
 f_k-f_{k+1}\ge\delta_k\min\{1/2,\kappa/(C_jD^2)\}
 >\min\{b_j/4,\kappa b_j/(2C_jD^2)\}.
\]
The phase budget $b_j$ permits at most
$5+4H_0D^2/\kappa+4H_1D^2b_j/\kappa$ such steps.
Sum with $\sum_jb_j<2U_0$, add lag steps, and apply the drop count.
Lemma~\ref{lem:awaybridge} ensures the actual stopping certificate has
passed by that point. No stronger assertion is made for general convex
functions: the extra geometric error bound matters.
\end{proof}
The geometric inequality, drop counting, and linear convergence of
adaptive AFW are established results. The candidate extension here is
the separated curvature dependence and its conversion to this particular
certificate, not a new away-step method or a priority claim.
The requirement $H_0>0$ in the earlier working version was stronger
than needed: the proof only needs positive caps at positive source
errors, which $H_1>0$ also ensures. For example,
$f(x)=t+(\max\{t,0\})^3$, $t=1/2-x$, on $[-1/2,1/2]$ has $H_0=0,H_1=6$,
zero strong-convexity modulus, and a positive selected-gap geometric
constant. The all-affine case $H_0=H_1=0$ is covered by the convex
backtracking bound, without invoking these caps.
For comparison, substituting
$L_Q=H_0+H_1\max_Q(f-f_Q^*)$ in a global-smoothness linear rate leaves
the entire global curvature multiplying the accuracy logarithm.
Here the explicit $H_1$ phase cost is proportional to $U_0$, with further
dependence inside $J_*$ through the certificate bridge. The new bridge
uses no $\mu$ or initial-level curvature $H_0+H_1U_0$; those quantities
are unnecessary for converting primal error into this stopping criterion.
This comparison
does not rule out deriving an equivalent specialization from a prior
adaptive AFW proof; its priority remains unresolved.
